% 本文件由 LaTeX 排版 Agent 根据有效 translation.json 自动生成。
% author=Augustine of Hippo; work_id=1310; title=论圣童贞

\chapter{论圣童贞}

% structure: p0001 source=h1 target=omitted reason=page-title-already-rendered-by-parent

论童贞\par

\EditorialNote{出自\WorkTitle{修正篇}2:23：在我写了\WorkTitle{论婚姻之善}之后，大家期待我写论圣童贞；我没有拖延：并且它是天主的恩赐，以及多么大的恩赐，以及要用何等谦逊来守护，尽我所能在一卷中阐述。本书以此开始，等等。}\par

1. 我们最近出版了一本名为\WorkTitle{论婚姻之善}的书，在其中我们也曾劝诫并现在劝诫\Person{基督}的贞女们，不要因为她们所领受的那更大恩赐，就看不起与自己相比的天主子民的父母；并且不要以为那些人（宗徒以橄榄树来比喻他们，使被接上的野橄榄不致骄傲），他们甚至借生育儿子来服侍将要来临的\Person{基督}，就因此功绩较小，因为按神圣权利，节制优于婚姻生活，虔敬的童贞优于婚姻。诚然，在他们身上，未来的事正在预备并产生出来，这些事现在我们以奇妙而有效的方式看到应验了，他们的婚姻生活也是预言的：因此，不是按人类愿望和欢乐的通常习惯，而是按天主极深的上智安排，在他们中某些人的多产获得了荣耀，某些人的不生育也变成了多产。但此时，对于那些被告知\BibleQuote{他们若不能自制，就让他们结婚}的人，我们必须使用不是安慰而是劝勉。但对于那些被告知\BibleQuote{谁能接受，就让他接受}的人，我们必须劝勉，使他们不要惊慌；并警告他们不要自高自大。因此，童贞不仅应当被表彰，使人爱慕，还应当被告诫，使人不致骄傲。\par

2. 这在我们的当前论述中已承担：愿\Person{基督}帮助我们，祂是童贞女的圣子，也是贞女的净配，按肉体生于童贞的母胎，按圣神在童贞婚姻中成婚。因此，既然整个教会本身是一位贞女，许配于一位丈夫\Person{基督}，正如宗徒所说，其成员何等堪受荣耀，他们在肉体本身中也保存了全教会以信德所保存的贞洁？这童贞模仿了她丈夫和主的母亲。因为教会也是一位母亲和一位贞女。因为如果她不是贞女，我们为谁的贞洁纯洁而忧虑？如果她不是母亲，我们称谁为子女？\Person{玛利亚}按肉体生育了这身体的元首，教会按圣神生育了那身体的肢体。在两者中，童贞不妨碍多产；在两者中，多产不除去童贞。所以，既然整个教会在身体和灵魂上都是圣的，然而整个教会并非在肉体上是贞女而是在圣神上是贞女，那么在那些在肉体上和圣神上都是贞女的肢体中，是何等更圣呢？\par

3. 在福音中记载了\Person{基督}的母亲和兄弟，即祂按肉体的亲属，当人传话给祂，他们站在外面，因为不能因人群到祂跟前时，祂回答说：\BibleQuote{谁是我的母亲？谁是我的兄弟？然后祂伸出手指着门徒说，这些人就是我的兄弟；凡遵行我圣父旨意的人，他就是我的兄弟、姊妹和母亲。}这除了教导我们：宁要按圣神的出身，也不要按肉体的亲属，并且人蒙福不是因为与义人和圣人有血肉之亲，而是因为藉着服从和追随，他们依附于他们的道理和行为。因此，\Person{玛利亚}因接受\Person{基督}的信德而更蒙福，胜于因孕育\Person{基督}的肉体。因为当一个人说：\BibleQuote{怀你的胎是有福的！}祂亲自回答说：\BibleQuote{是的，但更有福的是听天主的话而遵守的人。}最后，对于祂的兄弟，即按肉体的亲属，他们不信祂，那亲属关系有何益处？同样，\Person{玛利亚}作为母亲的血亲对她也无益，如果她没有以一种比在肉体中更蒙福的方式在心中怀有\Person{基督}。\par

4. 她的童贞本身也因此更蒙悦纳，因为\Person{基督}在她内成孕并非从将要侵犯的丈夫手中救回童贞，由祂自己保存；而是祂尚未成孕之前，就选择了她，因为她已经将自己奉献给天主，作为祂降生的来源。这一点由\Person{玛利亚}回答天使宣告她受孕的话显示出来：\BibleQuote{怎样这事能成就呢？因为我不认识男人。}这话她肯定不会说，除非她以前已将自己作为贞女奉献给天主。但由于以色列人的习惯尚不接受这一点，她许配给了一个义人，他不会用暴力夺去，反而会保护她免受暴力侵害她已奉献的童贞。尽管，即使她只说\BibleQuote{这事怎么能成就呢？}而没有加\BibleQuote{因为我不认识男人}，她当然也不会问一个女性如何生育所应许的儿子，如果她已打算行房事。她本可被命令继续作贞女，以便在她身上借着合宜的奇迹，天主子取得仆人的形状；但为了作为圣贞女们的榜样，免得有人认为只有她需要成为贞女，因为她获得了甚至无性交而怀孕的恩赐，她在还不知道自己将怀孕什么时，就将自己的童贞奉献给天主，以便在尘世和可朽的身体上模仿天上生活出于誓愿，而非命令；由于选择的爱情，而非服事的必要性。因此，\Person{基督}由一位贞女所生，而她在知道将从她所生者是谁之前，已决心继续保持童贞，祂宁愿认可而非命令圣童贞。这样，即使在祂取了仆人形状的这女性身上，祂也愿意童贞是自由的。\par

5. 因此，天主的贞女们没有理由忧伤，因为她们虽守贞，却也不能成为肉身之母。因为只有那一位，祂的诞生无与伦比，才适合由贞女生育。然而，童贞玛利亚的诞生是所有圣洁贞女的荣耀；她们与玛利亚一起，若遵行圣父的旨意，便是基督的母亲。因为按照上述祂的话，玛利亚也是基督的母亲，更值得赞美和祝福。\BibleQuote{凡遵行我在天之父旨意的人，就是我的弟兄、姊妹和母亲。}祂以属神的方式，在祂所救赎的子民中，彰显所有这些与祂的亲属关系：祂有圣洁的男女作为弟兄姊妹，因为他们都是天上产业的共同继承人。祂的母亲是整个教会，因为教会确实产生祂的肢体，即祂的信众。同样，每个虔诚的灵魂，若以最丰盛的爱德遵行圣父的旨意，在那些她为之劳苦的人中，直到基督在他们内成形，她也是祂的母亲。因此，玛利亚按肉身遵行天主的旨意，只是基督的母亲；但按圣神，她既是祂的姊妹，也是祂的母亲。\par

6. 正因如此，那一位女性，不仅按圣神，而且按肉身，既是母亲又是贞女。她按圣神确实是母亲，但不是我们元首的母亲，元首即是救主本身，她反而是由祂按圣神所生；因为凡信祂的人，其中也包括她，都理应被称为\BibleQuote{新郎的子女}；但她显然是祂肢体（即我们）的母亲，因为她以爱德合作，使信徒在教会中出生，成为那元首的肢体；但按肉身，她是元首本身的母亲。因为我们的元首，为了一桩伟大的奇迹，理应由一位贞女按肉身诞生，以此指明祂的肢体将按圣神，由教会这贞女诞生。因此，唯有玛利亚，按圣神和肉身，既是母亲又是贞女：既是基督的母亲，又是基督的贞女；而教会，在那些将要承受天主国度的圣人们中，按圣神，确实是基督的母亲，也完全是基督的贞女；但按肉身，并非完全如此，而是在某些人身上是基督的贞女，在某些人身上是母亲，但不是基督的母亲。因为已婚的信女和奉献给天主的贞女，因圣善的品行、出自纯洁之心的爱德、良好的良心和无伪的信德，因她们遵行圣父的旨意，按属神的意义都是基督的母亲。但那些在婚姻生活中按肉身生育子女的，所生的不是基督，而是亚当；因此她们努力，使她们的子女浸染于祂的圣事中，成为基督的肢体，因为她们知道她们所生的是什么。\par

7. 我说这番话，免得已婚的多产胆敢与贞女的贞洁竞争，并举出玛利亚本人为例，向天主的贞女说：\Quote{她按肉身拥有两件值得尊敬的事：贞洁和多产，因她既保持了童贞，又生了子。这福气，既然我们不能两者兼得，便分占了：你们作贞女，我们作母亲；你们在子女上所欠缺的，让你们所持守的贞洁作为安慰；对我们，得子的收益弥补我们失去的贞洁。}已婚信女的这番话，若她们真按肉身生出基督徒来，或许还可忍受；这样，除了贞洁，玛利亚肉身的生育能力之所以更卓越，仅仅在于她生了这些肢体的元首本身，而她们生了那元首的肢体。但现在，尽管那些为此而结婚并与丈夫同房、且只为得子，且对儿子别无他想，只想为他们赢得基督，并尽快这样做的人，用这种话争论；然而，基督徒并非由她们的肉身所生，而是后来才成为的：教会生出他们，因她以属神的方式是基督肢体的母亲，而教会也以属神的方式是基督的贞女。对于这神圣的生育，那些按肉身没有生出基督徒的母亲们，也同工合作，使他们成为她们知道自己按肉身不能生出的。然而，她们通过这一点同工合作，即她们自己也是贞女和基督的母亲，也就是说，在于\BibleQuote{以爱德行事的信德}。\par

8. 因此，任何肉身的生育都不能与圣洁的贞洁相比，即使是肉身的贞洁。因为贞洁本身不是因其为贞洁而受尊崇，而是因其已奉献给天主；虽然它在肉身内持守，却是因圣神的虔诚与热心而持守。由此，肉身的贞洁也是属神的，这是虔敬的节操所誓愿并持守的。因为，正如没有人会不洁地使用身体，除非那罪已在灵里怀胎；同样，没有人能在身体上保持端谨，除非贞洁已先栽植在灵里。此外，如果婚姻生活中的端谨，虽然在肉身内守护，却仍归功于灵魂，而非肉身——在灵魂的治理与引导下，肉身本身除了与婚姻的正当身份之外，不与任何人交合——那么，何况带着更大的尊荣，我们岂不更应将那种节操列入灵魂的美善中，即藉此将肉身的贞洁纯正为灵魂与肉身的创造者本身而誓愿、祝圣并持守的呢？\par

9. 是故，我们也不该相信，那些现今在婚姻中除了儿女别无他求，并将儿女归给基督的人，其肉身的多产能与童贞的丧失相提并论。诚然，在古代，为那按肉身将要来临的基督，在某个广大而具有预言性的民族中，肉身的种族本身是必要的；但如今，当基督的肢体可以从各种族、万民中被聚集到天主的子民和天国之城的时候，谁能接受神圣的童贞，就让他接受吧；唯有那不能节制的，才可结婚。因为，倘若某个富有的妇人为这项善工花费大量金钱，从不同民族买来奴隶使他们成为基督徒，她岂不以一种比任何（无论多么巨大的）子宫的多产都更为丰盛和众多方式，为基督肢体的诞生提供了准备？然而，她仍不敢将她的金钱与神圣童贞的奉献相提并论。但是，若为使将要出生的人成为基督徒，肉身的多产便有正当理由用来对抗贞洁的丧失，那么，若以巨大的金钱代价失去童贞，从而可以购买更多孩子使之成为基督徒，这比单从一个（无论如何多产的）子宫所能生的还要多产，那么此事将更为多产。但是，若说这话是极其愚蠢的，那么，已婚的信女们就当拥有她们自己的善，我们已在另一卷书中按似乎合宜的方式论述了这善；而她们更当（正如她们所惯常正确做的）在神圣的童贞女身上，尊崇那更好的善，即我们现下论述的善。\par

10. 因为，即或在此事上，已婚之人也不应以生有童贞女来与节欲者的功德相比：这并非婚姻的善，而是本性的善；是天主如此安排的，即无论人类两性的何种性交，无论是正当而诚实的，还是卑劣而不法的，所生的女性无一不是童贞女，却无一人生来便是神圣的童贞女；于是便发生这样的事：即使是出于淫乱，也能生出一个童贞女，但神圣的童贞女，却连婚姻也无法生出。\par

11. 我们自己也不是在童贞女身上提出她们是童贞女，而是提出她们是因虔诚的节欲而奉献于天主的童贞女。因为，我并非轻率地说，在我看来，已婚的妇人似乎比将要出嫁的童贞女更有福：因为前者拥有后者尚且渴望的，尤其若后者尚未与任何人订婚。前者学习取悦于她已许配的那一位；后者则要在众多人中取悦于她将许配的那一位（尚不知是谁）：她仅凭这一点便从人群中守护思想的贞洁，即她在人群中寻求的不是一个奸夫，而是一个丈夫。因此，那既不将自己置于众人面前以求爱慕（因为她要从众人中寻求一人的爱），也非现已找到此人便为己安排，思虑世俗之事，\Quote{怎样取悦丈夫}；而是如此爱上了\Quote{在世人中最美的}，以致她不能像玛利亚那样在肉身中怀有祂，却为自己在心灵中所怀的祂，也保守了肉身为童贞。这种童贞女，没有肉身的任何多产能生出她们：这不是血肉的产物。若要寻找她们的母亲，那就是教会。除了一位神圣的童贞女——她已许配给一位夫君基督，为被献上作贞洁的童贞女——之外，没有人能生下神圣的童贞女。她，并非完全在肉身，而是完全在精神上为童贞女，生下在肉身和精神上均为圣洁的童贞女。\par

12. 让婚姻拥有自己的善吧：不是在于生养儿子，而是在于以诚实、合法、端庄、以及团结的精神生养他们，并在他们出生后，以合作、有益的教育和认真的目的来养育他们；在于夫妻之间彼此保持床笫的信实；在于不破坏婚姻的圣事。然而，这一切都是人性的职责；但童贞的贞洁和因虔诚节欲而脱离一切性交的自由，却是天使的分内之事，是在可朽的肉身中实践永久的不可朽坏。让一切肉身的多产、一切婚姻生活的贞洁都屈服于此吧：前者不在人的能力之内，后者不在永恒之内；自由意志不拥有肉身的多产，天国不拥有婚姻生活的贞洁。确实，那些在肉身中已经拥有了某种不属于肉身之物的人，在那共同的永生中将拥有超乎他人的伟大之物。\par

13. 因此，那些认为这种节欲的善不是为了天国之故，而是为了现世之故才是必要的人，真是不可思议地缺乏智慧：仿佛，已婚之人因地上的牵挂而受困于越来越狭窄的各方面，而童贞和节欲之人却免于这种烦恼；仿佛仅仅为此原因，不结婚才更好，即可以逃避现时的困境，而非对未来生命有任何益处。而且，为使他们不显得是出于自己内心的虚妄而提出这虚妄的意见，他们引宗徒为证，那里说：\Quote{关于童贞女，我没有主的命令，但我既因主的仁慈而蒙信任，就提供意见。所以我认为，为了现时的需要，这为善是正确的，因为一个人这样是好的。}看，他们说，宗徒在此指明\Quote{为了现时的需要}，而不是为了未来的永恒。仿佛宗徒会顾及现时的需要，除非是为将来的事作准备和筹划；而他的一切作为，无不是指向永生。\par

14. 因此，我们必须避免的是目前的急需，但这种急需却妨碍了未来的某些美善；由于这种急需，婚姻生活被迫思虑世俗的事，如何取悦丈夫或妻子。这并不是说这些事使人远离天主的国，因为有些罪是借着命令而非劝诫来约束的，这是因为不遵守主的命令应受惩罚；但是，在天主国本身之内，若能更多思虑如何取悦天主，就能获得更大的福分，而由于婚姻的急需，这种思虑减少，福分也必然减少。因此他说：\BibleQuote{关于童贞女，我没有主的命令。}因为凡不遵守命令的，就是有罪的，应受惩罚。所以，既然娶妻或出嫁不是罪（如果是罪，就会以\Quote{命令}禁止），因此关于童贞女没有主的\Quote{命令}。但是，当我们避开或获得罪的赦免后，我们必须走向永生，那里有某种更卓越的光荣，不是赐给所有永生的人，而是赐给某些人；为了获得这光荣，仅仅从罪中获得释放是不够的，除非我们向释放我们的那一位许下某种不许可也不算过错、但许下并履行则是可称赞的事；他说：\BibleQuote{我赐予劝告，如同蒙主仁慈成为可信赖的人。}因为我不该吝啬忠信的劝告，我并非凭自己的功德，而是因天主的仁慈成为忠信的。\BibleQuote{所以，我因目前的急需认为这是好的。}他说，关于我没有主命令而只赐予劝告的事，即关于童贞女，我因目前的急需认为这是好的。因为我知道目前时日的急需，婚姻为此服务，它迫使人对天主的事思虑不足，不足以获得那并非所有人都有的光荣，尽管他们留在永生和救恩中：\BibleQuote{因为星辰与星辰的光辉不同；同样，死人复活也是如此。}因此，\BibleQuote{人这样存留是好的。}\par

15. 此后，同一宗徒接着说：\BibleQuote{你已与妻子结合，不要寻求解脱；你已脱离妻子，不要寻求妻子。}这两句话中，第一句属于命令，违反它就是违法的。因为除非因奸淫，不可休妻，正如主自己在福音中所说的。但他所加的\BibleQuote{你已脱离妻子，不要寻求妻子}是劝戒的话，而非命令；因此这样做是许可的，但不做更好。最后，他立刻补充说：\BibleQuote{你若娶妻，并没有犯罪；童贞女若出嫁，也没有犯罪。}但在他说了前面的话\BibleQuote{你已与妻子结合，不要寻求解脱}之后，他并没有加上\Quote{你若解脱了，也没有犯罪}，是吗？因为他之前已经说过：\BibleQuote{至于那些已婚的人，我命令——不是我，而是主——妻子不可离开丈夫；但如果她离开了，就该独身，或与丈夫和好；}因为可能她离开并非因自己的过错，而是丈夫的过错。然后他说：\BibleQuote{丈夫也不可休妻，}这是主的命令；他当时也没有加上\Quote{他若休了妻，也没有犯罪}。因为这是命令，不服从就是罪；而不是劝戒，你若不愿采用，将获得较少的美善，而非作恶。因此，在他说了\BibleQuote{你已脱离妻子，不要寻求妻子}之后，因为他不是下命令以免作恶，而是赐予劝戒以行更美的事，他立刻加上：\BibleQuote{你若娶妻，并没有犯罪；童贞女若出嫁，也没有犯罪。}\par

16. 然而，他接着说：\BibleQuote{但这样的人肉身将受苦难，而我愿意你们免受这些。}他这样劝勉童贞和持续节欲，也略略警告人远离婚姻，但完全本着谦逊，不是作为邪恶和非法的事，而是作为负担和麻烦的事。因为招致肉身的羞辱是一回事，而肉身受苦难是另一回事：前者是犯罪的行为，后者是受苦的劳作，即使为了最尊贵的职责，人们也多半不拒绝。但就婚姻而言，在现今这个时代，不再通过生育子女为即将来临的基督服务，却要承担宗徒所预告那些结婚的人将有的肉身苦难，那是极其愚蠢的，除非那些不能节欲的人害怕因撒殚的诱惑而陷入该受惩罚的罪。至于他说他愿意那些人免受苦难，即那些他预言将受肉身苦难的人，我暂时想不到更稳妥的解释，除非他不想明说并详细描述这种肉身苦难，即嫉妒婚姻生活的猜疑、生儿育女的操劳、无子女的恐惧和悲伤。因为有多少人，在将自己束缚于婚姻之锁后，不被这些情感所牵引和驱策呢？我们不应夸大这些，以免不给那些宗徒认为应得免去的人留情面。\par

17. 惟有藉此我简要陈明的事，读者应当提防那些人：他们对于所写的这段话，即\BibleQuote{但这样的人在肉身方面将受困难，而我却宽免你们，}，虚假地谴责婚姻，仿佛此句话间接谴责了婚姻；好像他不愿说出谴责的话，当他说\BibleQuote{但我宽免你们；}时，诚然，当他宽免他们时，他没有宽免自己的灵魂，反而虚假地说：\BibleQuote{而且，你若娶妻，你并没有犯罪；若童女出嫁，她也不犯罪。}凡相信或愿意相信圣经的人，他们好比为自己预备了说谎的自由，或辩护自己邪见的途径，无论何时他们持守与纯正道理相违的意见。因为若有人从神圣经卷中引述明确的文句来驳斥他们的错误，他们便以此为盾牌；藉此自卫以对抗真理，却暴露自己成为魔鬼的伤口：即声称本书的作者在这一点上没有说实话，有时是为了宽免软弱者，有时是为了恐吓轻蔑者；就像碰巧遇到一个场合，用以辩护自己的邪见一样；如此，既然他们宁愿辩护而不愿改正自己的见解，他们就试图破坏圣经的权威，唯独藉着圣经才能折断一切骄傲顽固的颈项。\par

18. 因此我劝勉追求永久节德和圣洁童贞的男女，他们应将自己的善置于婚姻之上，但不可判断婚姻为恶；要明白宗徒所说的话绝不是出于诡诈，而是出于明白的真理：\BibleQuote{凡结婚的，作得好；凡不结婚的，作得更好；而且，你若娶妻，你并没有犯罪；若童女出嫁，她也不犯罪；}稍后又说：\BibleQuote{但她若这样坚持，按照我的意见，更是有福的。}为使这意见不被视为人的意见，他加上：\BibleQuote{可是我认为我也有天主的圣神。}这是主的教训，这是宗徒的教训，这是真实的，这是健全的：如此选择更大的恩赐，以致较小的不被谴责。天主的真理，在天主的圣经中，胜于任何人心中或肉身上的童贞。愿纯洁被珍爱，以致真理不被否认。因为那些相信宗徒的口舌，在恰恰是赞许身体童贞之处，竟未被说谎的腐败所玷污，他们对自己肉身还能有怎样的邪恶想法呢？因此，首先且最重要的是，凡选择童贞之善的人，要最坚定地持守圣经绝无虚言；并且因此，所说的这句话也是真实的：\BibleQuote{你若娶妻，你并没有犯罪；若童女出嫁，她也不犯罪。}并且让他们不要以为，若婚姻不是罪恶，童贞洁德如此伟大的善就被贬低了。反之，她应当因此确信，为她预备了更大光荣的棕榈枝，因为她不怕若结婚会被定罪，而渴望因未结婚而接受更尊贵的荣冠。因此凡愿意不结婚而生活的人，不要逃避婚姻如同罪的陷阱；而要超越它，如同较小的善的山丘，以便他们能安息在更大的善——节德——的山上。诚然，正是以此为条件，这座山丘才有人居住：人不能随意离开它。因为\BibleQuote{女人当丈夫活着的时候，是被束缚的。}然而，人从它如同从台阶上，升至寡居的节德；但为了童贞的节德，人必须或者藉不接受求婚者而回避它，或者藉抢先于求婚者而越过它。\par

19. 但恐怕有人以为两种善工，即善的与更好的，报酬会相等，为此必须反驳那些人：他们如此解释宗徒的那句话\BibleQuote{但我以为，由于当前的困难，这是好的}，竟说童贞有用不是为了天国，而是为了现世；仿佛在永生中，选择这更好部分的人不比其余的人多什么。在这讨论中，当我们论及同一宗徒的那句话\BibleQuote{但这样的人在肉身方面将受困难，而我却宽免你们；}时，我们遇到了另一些争辩者，他们远不将婚姻等同于永久的童贞，却完全谴责它。因为两者都是错误：或将婚姻等同于圣洁的童贞，或谴责它；由于彼此过度偏离，这两种错误公开碰撞，因为它们不愿持守真理的中道：藉此，藉着确定的理性和圣经的权威，我们既发现婚姻不是罪，又不将它等同于童贞或甚至寡居的洁德之善。有些人诚然因追求童贞，而视婚姻如同通奸一样可憎；但另一些人因辩护婚姻，竟要永久节德之卓越不配得比已婚洁德更多；仿佛\Person{苏撒纳}的善是贬抑\Person{童贞玛利亚}，或\Person{童贞玛利亚}的更大的善应当谴责\Person{苏撒纳}。\par

20. 因此，愿宗徒对已婚者或将要结婚的人所说的话，\BibleQuote{但我宽免你们}，并不意味着他不愿说出已婚者在来世应受的惩罚。愿那由达尼尔从现世审判中解救出来的女子，不被保禄投入地狱！愿她在基督审判台前，因坚守对丈夫的忠贞而选择面对危险或死亡的夫妻贞洁，不成为她的惩罚！那句话说：\BibleQuote{我宁愿落在你们手中，也不愿得罪天主}，其意义何在？如果天主并非要因她坚守夫妻贞洁而解救她，而是要因她结婚而惩罚她，那又该如何？如今，每当神圣圣经的真理为夫妻贞洁辩护，反对那些诽谤和控告婚姻的人时，圣神就不断为苏撒纳辩护，使她免受假见证的控告，并更加费力地解救她脱离虚妄的指控。因为那时只针对一个已婚妇女，现在却针对所有已婚者；那时是隐秘的、不实的通奸指控，现在却是真实公开的婚姻指控。那时一个妇女因不义长老的诬告而受审，现在所有丈夫和妻子却因宗徒不愿说的话而被控告。他们说：正是对你们的定罪，他在说\BibleQuote{但我宽免你们}时保持沉默。这话是谁说的？当然是那位在上文说过\BibleQuote{你若娶妻，没有犯罪；童女若出嫁，也没有犯罪}的人。那么，为什么你们要在他因谦逊而保持沉默之处，怀疑他对婚姻的指控，而在他公开说话之处，却认不出他对婚姻的辩护呢？难道他用沉默定他们的罪，而用言语宣告他们无罪吗？现在，控告苏撒纳的不是婚姻，而是通奸本身，这难道不比控告宗徒的教导是谎言更为温和吗？在如此巨大的危险中，我们还能做什么，除非确定无疑地认为，贞洁的婚姻不应被定罪，正如神圣圣经绝不撒谎一样确定无疑？\par

21. 这里有人会说：这一切与圣洁的童贞或永久贞洁有何关系？本次演讲正是为了阐述这些。我首先要回答的是，如上所述，那更大之善的光荣之所以更大，是因为为了获得它，人们克服了婚姻生活的善，而非避免了婚姻的罪。否则，永久贞洁就足以不被特别赞扬，而仅不受责备——如果维护它是因结婚被视为罪恶的话。其次，既然如此卓越的恩赐必须藉神圣圣经的权威，而非人的判断来劝勉人，我们就不能以平庸或附带的方式辩护，以免神圣圣经本身在任何事上显得说谎。因为那些通过定婚姻的罪来强迫童贞者继续守贞的人，与其说是劝勉，不如说是吓阻他们。因为如果他们认为圣经中紧邻上文写着的\BibleQuote{凡出嫁自己的童女，是做得好}是假的，他们怎能确信所写的\BibleQuote{凡不出嫁的，做得更好}是真的呢？但是，如果他们对圣经所言婚姻之善（这由同一最真实的神谕权威所证明）毫无疑问地相信，他们就会以炽热而自信的热忱，奔向自己更美好的部分。因此，对于我们承担的事业，我们已经说得足够，并尽己所能地表明：无论是宗徒说的\BibleQuote{我想为现世的急需，这样是好的}，不应理解为圣洁的童贞在此生比已婚的信女更好，而在天国和来世平等；还是另一处论及结婚者所说的\BibleQuote{这样的人肉身必有苦难，但我宽免你们}，不应理解为他不愿说出婚姻的罪和定罪。事实上，由于误解，两个彼此相反的错误分别附着于这两句话：关于现世急需的那句话，被那些主张结婚者与不结婚者平等的人按自己的意思解释；而\BibleQuote{但我宽免你们}这句话，则被那些妄自定罪婚姻的人利用。但按照圣经的信德和健全教义，我们既说婚姻无罪，又将其善置于不仅低于童贞贞洁，也低于寡居的贞洁之下；并说已婚者现世的急需阻碍了他们的功德——不是阻碍永生，而是阻碍那为永久贞洁所保留的卓越光荣和尊荣；并且，现今除非不能节欲的人，结婚并无益处；至于肉身的苦难（这来自肉体的情欲，没有它，不节欲者的婚姻无法存在），宗徒既不愿缄口不言——因为那是真实的警告，也不愿更详尽地展开——因为体恤人性的软弱。\par

22. 如今，藉着神圣圣经最清晰的见证——按我们记忆所及的小小度量所能忆起的那些——愿它更清楚地显明：我们之所以选择永久的贞洁，并非为了这现世的生命，而是为了那在天国中所应许的未来生命。然而，谁不会在同一位宗徒稍后所说的话中注意到这一点呢？他说：\BibleQuote{没有妻子的，所挂虑的是主的事，想怎样悦乐主；但娶了妻子的，所挂虑的是世俗的事，想怎样悦乐妻子。未出嫁的妇女和童贞女是分开的：未出嫁的妇女挂虑主的事，为求身心圣洁；但出嫁的妇女挂虑世俗的事，想怎样悦乐丈夫。}显然，他并没有说，她所挂虑的是在这世上无忧无虑的状态，为的是度过没有更沉重烦恼的时光；他也没有说，未出嫁的妇女和童贞女是分开的——即与那出嫁的妇女区别和分离——是为了让未出嫁的妇女在今生无忧无虑，以避免世间的烦恼——这是已婚妇女所不能免的。而是说：\BibleQuote{她所挂虑的是主的事，想怎样悦乐主；并挂虑主的事，为求身心圣洁。}除非有人愚昧好辩到如此地步，竟敢断言：我们之所以想\BibleQuote{悦乐主}，不是为了天国，而是为了现世；或者说，她们之\BibleQuote{身心圣洁}，是为了今生，而非为了永生。相信这一点，岂不成了比众人更可怜的人吗？因为宗徒说：\BibleQuote{如果我们只在今生寄望于基督，我们就比众人更可怜了。}怎么？那将自己饼分给饥饿之人的人，若只是为了今生而行，便是愚昧；而那克制自己身体以至贞洁、以致于连婚姻的交接也放弃的人，若在天国中毫无所获，难道就能算为明智吗？\par

23. 最后，让我们听主自己对此事所发表的最明确的判决。因为，当祂以神圣而可畏的方式论及夫妻不可分离——除非因奸淫——之后，祂的门徒对祂说：\BibleQuote{如果同妻子的情形是这样，倒不如不结婚了。}祂对他们说：\BibleQuote{这话不是人人都能领受的。因为有些阉人是从母胎生来的，有些是被人阉的，也有些是自阉的——他们为了天国而自阉。谁能领受，就领受吧。}有什么比这更真实、更清楚的呢？基督——那真理、天主的德能和智慧——说：那些出于虔诚之目的而节制自己、不娶妻子的人，是为了天国而自阉。而人的虚荣却以不敬的鲁莽反对这一点，声称这样做的人只是为了躲避婚姻生活的眼前烦恼，但在天国中并不比别人更有份。\par

24. 但是，天主藉先知依撒意亚论及阉人所说的，祂要在祂的屋内和墙内赐给他们一个名号的地方，远胜于子女，这除了指那些为天国而自阉的人之外，还能指谁呢？因为对于那些身体器官无力以致不能生育的人（如富人和国王的阉人），当他们成为基督徒并遵守天主的诫命时，如果他们有这个意图——若能娶妻，他们就会娶——那么他们在天主的屋内，与那些合法且贞洁地生育子女并教导儿子们寄望于天主的已婚信友同列，这确实足够了；但他们不会得到一个比子女更\Emph{好}的名号。因为不娶妻并非出于灵魂的德行，而是出于肉身的需要。任凭谁愿意争辩说先知预言的是那些遭受身体阉割的阉人；那也有助于我们承担的论题。因为天主并没有将这些阉人置于那些在祂屋内没有位置的人之上，而是显然置于那些在生育子女上保持婚姻功绩的人之上。因为当祂说：\BibleQuote{我要赐给他们一个更好的地方}，这表明已婚者也得着一个地方，但远逊于前者。因此，允许在旧约以色列中没有的上述肉身的阉人将在天主的屋内；因为我们看到他们本身，虽不成为犹太人，却成为基督徒；而先知所说的不是那些因节制志向而寻求不结婚，为天国而自阉的人：难道有人如此疯狂地反对真理，以致相信那些在肉身中成为阉人的人在天主屋内比已婚者有更好的位置，并主张那些虔敬守贞、克己肉身甚至轻视婚姻、不是在身体上而是在贪欲的根上自阉、在尘世必朽状态中实践天上和天使般生活的人，其功绩与已婚者相等；并且身为基督徒，却反驳基督赞美那些不是为现世而是为天国而自阉的人，声称这有益于现世生命而非未来？这些人除了断言天国本身属于我们现在所处的现世生命之外，还能有什么别的呢？因为盲目的傲慢为什么不应该甚至走向这种疯狂？还有什么比这种断言更疯狂的？因为虽然有时教会——即使是现在的教会——被称为天国，当然这样称呼是为了这个目的，因为它正在为未来的永生而聚集。因此，尽管它既有今世生命的应许，也有来世生命的应许，但在它的一切善工中，它不注目\BibleQuote{看得见的事，而是看不见的事；因为看得见的是暂时的，看不见的是永远的}。\newline{}\par

25. 圣神也确实没有不发表明确而不可动摇的言论来对抗这些极其无耻和疯狂顽固的人，并用不可攻破的防御从羊圈中击退他们的攻击，如同野兽一般。因为论到阉人，祂说了\BibleQuote{我要在我的屋内和我的墙内赐给他们一个名号的地方，远胜于子女}之后，免得任何过于属肉体的人以为这些话语中有任何暂时的事可以指望，祂立刻加上：\BibleQuote{我要赐给他们一个永远的名号，永不消逝}；好像祂说：你为什么退缩，不敬的盲目？你为什么退缩？你为什么将你悖逆的云雾倾倒在真理的晴空之上？为何在如此明亮的圣经之光中，你仍从黑暗中寻求陷阱？为何你只将暂时的利益许诺给操练节制的圣人？\BibleQuote{我要赐给他们一个永远的名号}；为什么当人们远离一切性交，并且正因他们戒除这些，思念主的事，怎样能悦乐主时，你却试图将他们引向现世的利益？\BibleQuote{我要赐给他们一个永远的名号}。为什么你争辩说圣洁的阉人为之自阉的天国只能在这生命中被理解？\BibleQuote{我要赐给他们一个永远的名号}。即使你在这里试图把\Quote{永远}这个词理解为持续很长时间，我还要加上，堆积，踏进去：\BibleQuote{永不消逝}。你还寻求什么？你还说什么？这个永远的名号，无论它是什么，是赐给天主的阉人的，它无疑标志着某种特殊而卓越的荣耀，不会与许多人共享，尽管他们处于同一王国和同一屋内。也许正因此，它被称为名号，以将那些领受它的人与其他的人区分开来。\newline{}\par

26. 那么，他们说，在葡萄园工作结束时，同等付给所有人的那一个银钱是什么意思呢？无论是给从第一小时工作的人，还是给只工作了一小时的人？它无疑表示某种所有人都将共同拥有的东西，例如永生本身、天国本身，凡天主所预拣、召叫、成义、光荣化的人，都将在此。\BibleQuote{因为这必朽坏的必须穿上不朽坏的，这必死的必须穿上不死的。}这就是那银钱，所有人的工资。然而\BibleQuote{星的光荣不同，复活死人也是如此。}这些是圣人们不同的功德。因为，若那银钱象征天堂，难道不是所有星都共同在天上吗？然而，\BibleQuote{日有日的光荣，月有月的光荣，星有星的光荣。}若那银钱象征身体的健康，当我们健康时，难道不是所有肢体都共享健康吗？若这健康持续至死，难道不是在所有肢体中一样而且平等吗？然而，\BibleQuote{天主将肢体，每一个，安置在身体上，如祂所愿；}使得全体不成为一只眼，也不全是听觉，也不全是嗅觉：其余无论什么，都有其特性，尽管它与所有肢体一样健康。如此，因为永生本身将平等地给予众人，一个相等的银钱分配给了所有人；但是，因为在永生本身中功绩的光将以区别照耀，在圣父的家里有\BibleQuote{许多住处}；而且，借此，在不相异的银钱中，一人不比另一人活得更长；但在许多住处中，一人以更大的光荣受到尊荣。\par

27. 因此，天主的圣人们，男孩和女孩，男人和女人，未婚男子和女子，前进吧，坚持到底。更甜美地赞美主，你们更丰富地思念祂；更幸福地盼望祂，你们更热切地事奉祂；更热烈地爱祂，你们更殷勤地中悦祂。束上腰，点着灯，等候主从婚筵回来。你们要将一首新歌带到羔羊的婚筵，用你们的琴弹唱。那歌当然不是全地所唱的歌，关于全地曾说：\BibleQuote{你们要向主唱新歌；全地都要向主歌唱。}而是除你们以外无人能唱的歌。因为在默示录中，有一位比其他人更蒙羔羊喜爱的、曾靠在祂胸前的、惯于汲取并涌出天主圣言超越天上奇事的人，曾看见你们。他看见你们，十四万四千位圣洁的弹琴者，身体上纯洁的童贞，心里无玷的真理；他论到你们说，你们跟随羔羊，无论祂往哪里去。我们想，这羔羊往哪里去？在那里，除了你们，无人敢也不能跟随。我们想祂往哪里去？进入怎样的幽谷和草场？我想，那里有喜乐的青草；不是这世界的虚幻喜乐、疯狂的谎言；也不是在天主的国里为其余非童贞者所预备的喜乐；而是与众不同的喜乐份额。基督的童贞者的喜乐，是来自基督、在基督内、与基督一起、跟随基督、透过基督、为了基督的喜乐。基督的童贞者特有的喜乐，不同于那些虽属基督却是非童贞者的喜乐。因为不同的人有不同的喜乐，但没有一种喜乐像这样。去吧，进入这些，跟随羔羊，因为羔羊的肉体也确实是童贞的。因为祂在自己长大时保留了这一点，而祂的降孕和出生并未从祂的母亲那里取走这一点。你们以心灵和肉体的童贞，配得跟随祂，无论祂往哪里去。因为跟随不就是效法吗？因为\BibleQuote{基督为我们受了苦}，给我们留下了榜样，正如宗徒伯多禄所说，\BibleQuote{使我们跟随祂的足迹。}各人在他所效法的事上跟随祂：不是就祂是天主唯一圣子、万物由祂而造而言，而是就祂作为人子，在自己身上展示了我们应当效法的事而言。祂有许多事是给所有人效法的；但肉体的童贞不是给所有人的；因为那些已经不再是童贞的人，没有办法成为童贞。\par

28. 因此，其余失去童贞的信徒，可以跟随羔羊，不是无论祂往哪里去，而是按他们所能的程度跟随。他们处处都能跟随，除非当祂行走在童贞的恩宠中。\BibleQuote{神贫的人是有福的；}效法祂，祂\BibleQuote{本是富足的，却为你们成了贫穷的。}\BibleQuote{温良的人是有福的；}效法祂，祂说：\BibleQuote{你们跟我学，因为我是良善心谦的。}\BibleQuote{哀恸的人是有福的；}效法祂，祂\BibleQuote{曾为}耶路撒冷\BibleQuote{哭泣。}\BibleQuote{饥渴慕义的人是有福的；}效法祂，祂说：\BibleQuote{我的食物就是承行派遣我者的旨意。}\BibleQuote{怜悯人的人是有福的；}效法祂，祂帮助了那个被强盗打伤、半死躺在路上无望的人。\BibleQuote{心里洁净的人是有福的；}效法祂，祂\BibleQuote{没有犯罪，口中也没有诡诈。}\BibleQuote{缔造和平的人是有福的；}效法祂，祂为迫害自己的人说：\BibleQuote{父啊，宽恕他们，因为他们不知道自己在做什么。}\BibleQuote{为义而受迫害的人是有福的；}效法祂，祂\BibleQuote{为你们受苦，给你们留下榜样，叫你们跟随祂的足迹。}凡效法这些事的人，就在这些事上跟随羔羊。但即使是已婚的人也能走在这些足迹中，虽然不能完全踏在同一个脚印里，却走在同一条路上。\par

29. 但是，看，那只羔羊行走在一条童贞之路上，那些已经失去且无法挽回的人，怎能跟随祂呢？所以，你们，你们这些童贞者，要跟随祂；你们也要在那里跟随祂，因为唯独在此事上，无论祂往哪里去，你们都跟随祂：因为对于任何其他圣洁的恩赐，我们可以劝勉已婚者跟随祂，唯独此恩赐她们已失去而无法挽回。所以，你们要跟随祂，以恒心持守你们以热忱所誓许的。你们要尽力前行，免得童贞之善从你们那里消失，而你们无法使之复得。其余的众多信友将看见你们，他们不能在此事上跟随羔羊；他们将看见你们，却不嫉妒你们；因着与你们同乐，他们自身所没有的，将在你们身上拥有。因为他们不能唱那首也是属于你们的新歌；但他们能听见，并因你们如此卓越的美善而欢喜；但你们，既要唱又要听，因你们所说的正是你们自己听到的，将怀着更大的幸福欢跃，并怀着更大的喜乐为王。但他们不会因你们更大的喜乐而忧伤，尽管他们缺乏此喜乐。实在，那只羔羊，你们将跟随祂，无论祂往哪里去，祂不会抛弃那些不能跟随祂到你们所能到之处的人。我们所说的羔羊是全能者。祂既要在你们前面走，又不会离开他们，当天主成为万物中的万有时。那些拥有较少的人，不会因厌弃而离开你们；因为在那里没有嫉妒，差异和谐共存。所以，你们要鼓起勇气，要有信心，要刚强，要持守，你们这些向上主你们的天主誓许并偿还终身节欲之愿的人，不是为了现世，而是为了天国。\par

30. 你们这些尚未誓许此愿，且能领受的人，领受吧。以恒心奔跑，好能获得。你们各人要献上自己的祭品，进入上主的庭院，不是出于不得已，而是由于自己的意志自主。因为不能说\Quote{你们不可结婚}，如同说\BibleQuote{你们不可奸淫，不可杀人}那样。前者是必须履行的，后者是自愿奉献的。后者若实行了，便受赞美；前者若未遵守，便受谴责。在前者，主命令我们当尽的本分；但在后者，若你们花费了额外的，祂回来时会偿还你们。请想想，在祂的城墙内，\BibleQuote{一个名为……的地方，远远胜过儿女之名}（无论那是什么）。请想想那里的\BibleQuote{永恒之名}。谁来说说那名字是怎样的？但无论它是什么，它将是永恒的。因着相信、盼望和爱慕这一点，你们已能避开婚姻——并非因它被禁止，而是如允许之事般超越它。\par

31. 由此，这项奉献的伟大性——我们已按我们的力量劝勉你们承担——越是卓越和神圣，就越令我们的焦虑提醒我们，不仅要论及最荣光的贞洁，也要论及最安全的谦逊。当这些如此终身贞洁的誓愿者，将自己与已婚者比较时，发现按照圣经，后者无论在行动和报酬上，还是在誓愿和赏报上都较低劣时，那么，所记载的话应立即进入他们的思想：\BibleQuote{你越伟大，就越要在一切事上谦卑自己：这样你必在天主面前获得恩宠。}各人的谦逊尺度由其自身的伟大程度而定：骄傲对此充满危险，它对人越强大，就越设下更大的陷阱。随之而来的是嫉妒，如同它的女儿尾随其后；确实，骄傲立刻生出嫉妒，且嫉妒从不缺少这样的女儿和同伴。这两个邪恶，即骄傲和嫉妒，使魔鬼成为魔鬼。因此，针对骄傲——嫉妒之母——整个基督徒纪律主要与之作战。因为基督徒纪律教导谦逊，由此获得并保持爱德。关于爱德，说了\BibleQuote{爱德不嫉妒}之后，如同我们询问原因，它如何不嫉妒，立刻又加上\BibleQuote{不自大}，好像说：它不嫉妒，正是因为它也没有骄傲。因此，基督，谦逊的导师，首先\BibleQuote{空虚了自己，取了奴仆的形体，成了人的模样；且以人的形态出现，祂谦卑了自己，服从至死，且死在十字架上。}但祂的教导本身，如何谨慎地建议谦逊，如何热切而迫切地命令谦逊，谁能轻易阐述并汇集所有证据来证明此事呢？凡愿意者，可以尝试做或做这事，即撰写一篇关于谦逊的专论；但本作品的目标不同，且所讨论的题目如此重大，以至于主要须防范骄傲。\par

32. 为此，主屈尊提示我记起关于谦逊的基督训导中的几个见证，我将提出来，或许已足敷我用。祂初次向门徒所作的较长的讲论，便是从此开始的。\BibleQuote{神贫的人是有福的，因为天国是他们的。}我们毫不争议地将这些视为谦卑的人。祂特别称赞那位百夫长的信德，说在以色列中未曾见过如此大的信德，因为他以极大的谦卑相信，以至说：\BibleQuote{我不堪当你到舍下来。}为此，玛窦也说他没有别的理由\BibleQuote{来到}耶稣跟前（而路加最清楚地表明他并未亲自来到祂面前，而是派了他的朋友去），只是由于他最忠信的谦卑，他本人比他所派的人更来到祂面前。为此，先知也说：\BibleQuote{上主至高，却垂顾卑微的人；高傲的人，祂只远远认识；}无疑，因他们不来到祂面前。为此，祂对那客纳罕妇人也说：\BibleQuote{妇人，你的信德真大！就照你的愿望成就吧！}祂此前曾称她为狗，并答复说儿女的饼不可扔给她。她以谦卑接受这话，说：\BibleQuote{主，是啊！可是小狗也吃主人桌子上掉下来的碎屑。}这样，她藉着不断的呼喊未得到的，却藉着谦卑的承认赢得了。为此，也描述了两个人上圣殿祈祷，一个是法利塞人，另一个是税吏，为了那些自以为义而轻视他人的人，罪的承认被置于功绩的列举之前。法利塞人确实因那些使他大为自满的事而感谢天主。他说：\BibleQuote{我感谢你，因为我不像其他的人，勒索、不义、奸淫，也不像这个税吏。我每周两次禁食，我献上一切所得的十分之一。}但是税吏远远站着，连举目望天都不敢，只是捶着胸膛说：\BibleQuote{天主，可怜我这个罪人吧！}接着是神圣的判决：\BibleQuote{我实在告诉你们：这税吏下去，到他家里，成了正义的，而那个法利塞人却不然。}然后表明了为什么如此公正：\BibleQuote{因为凡高举自己的，必被贬抑；凡贬抑自己的，必被高举。}因此，可能发生这样的事：一个人既逃避真实的恶，也反省自己身上的真实善，并为这些善而感谢\BibleQuote{光之父，一切美好的赠与和一切完善的恩赐都是从祂那里降下来的}，却因傲慢之罪而被拒绝，如果因骄傲，即使只在思想中（在天主面前），他侮辱了其他罪人，特别是在祈祷中承认他们的罪时，对于这些人，应当给予的不是傲慢的责备，而是不绝望的怜悯。当祂的门徒彼此争论谁最大时，祂将一个小孩子放在他们眼前，说：\BibleQuote{你们若不变成如同小孩一样，决不能进天国。}祂不是主要推崇谦逊，并将伟大的功劳置于其中吗？或者，当载伯德的儿子们渴望坐在祂身边的高位上，祂如此回答，使他们宁可想到要喝祂苦难的杯（祂在其中贬抑自己，听命至死，且死在十字架上），而不是以骄傲的欲望要求优先于他人时，祂表明了什么呢？岂不是表明祂要赐予他们高举，只要他们首先跟随祂这位谦逊的老师？而现在，当祂即将去受难时，祂洗了门徒的脚，并最明确地教导他们为同门和同仆做这事，就是祂这位主和老师为他们所做的，祂是多么推崇谦逊啊！为了推崇谦逊，祂还选择了那个时刻，那时他们正迫切地注视祂，仿佛祂即将死去；无疑，他们将特别记住这一点，即他们要效法的老师最后指给他们的。但祂在那时做了这事，当然祂也可以在从前与祂们相处的其他日子做这事；如果那时做了，这训导固然会传下来，但肯定不会被如此接受。\par

33. 既然所有基督徒都必须守护谦逊，因为他们从基督得名基督徒，而任何人若仔细研读祂的福音，就会发现祂是谦逊的老师；那么，那些在任何大善中超越其余人的人，尤其应当成为此德行的追随者和守护者，以便他们能谨慎持守我开头所录的那句话：\BibleQuote{你越大，越要在一切事上谦卑自己，这样你才能在天主面前获得恩宠。}为此，由于持久的节德，尤其是童贞，是天主圣人的大善，他们必须万分警惕，免得它被骄傲所败坏。\par

34. 宗徒\Person{保禄}谴责邪恶的未婚女子，她们好奇而多言，并说这缺点来自闲散。他说：\BibleQuote{但同時，她們因閑散，便學習到處串門，不僅閑散，而且好奇多言，說不當說的話。}關於這些，他先前曾說：\BibleQuote{但不要收年輕的寡婦，因為當她們在享樂中度過時光後，就願意在基督內出嫁；她們受到譴責，因為她們背棄了當初的信德。}然而他不是說她們出嫁，而是說\BibleQuote{她們願意出嫁。}因為她們中的許多，不是因為對崇高目標的愛，而是因為害怕公開的羞辱，才被勸阻而不出嫁，而這羞辱本身也來自驕傲，因為人懼怕得罪人過於得罪天主。因此，這些願意出嫁卻因不能免受懲罰而不出嫁的人——她們與其被慾火的隱秘火焰在良心上燒毀，還不如出嫁；她們對自己的宣願後悔，覺得自己的宣認厭煩——除非她們改正並端正內心，再次以敬畏天主克服慾望，就必須被算在死人之中；不論她們是在享樂中度日——為此宗徒說：\BibleQuote{但那享樂度日的寡婦，雖活著也是死的。}——還是在勞苦和齋戒中度日，這若沒有內心的改正，便是無用的，只為炫耀而非改善。我本人，對於這樣的人——在其身上驕傲本身也混亂，因良心的傷口而流血——不要求大的謙遜。對於那些醉酒者、貪婪者，或任何其他該受罰的疾病——同時他們擁有肉體節制的宣願，但由於邪惡的品行而與自己的名聲不符——我不要求他們對虔誠的謙遜如此關切；除非這些邪惡者甚至敢於炫耀自己，他們還嫌這些懲罰來得太晚。我也不是在討論那些有某種取悅意圖的人——不論是穿著比這如此偉大宣願所需的更優雅的服裝，還是用顯著的束頭方式，或是使頭髮鼓起，或是用如此柔軟的頭巾，以致下面的細網顯露出來；對這些人，我們必須給予訓誡，還不是關於謙遜，而是關於潔德本身，或童貞的端庄。請給我一個宣願永久節制、且免受這些及所有類似缺陷和行為污點的人；我為這一個擔心驕傲，我為這如此大的善而害怕傲慢的膨脹。一人身上值得自滿的越多，我就越擔心，恐怕他因自滿，而不中悅那\BibleQuote{拒絕驕傲的人，卻賜恩寵給謙遜的人}。\par

35. 當然，我們要在基督自己身上默想童貞純潔的首要教訓和模範。那麼，關於謙遜，我還能給守貞者什麼訓誡，比祂對眾人說的更進一步呢？祂說：\BibleQuote{你們跟我學吧，因為我是良善心謙的。}之前祂提及自己的偉大，並且願意表明這件事本身——祂是何等偉大，又為我們成了何等微小——祂說：\BibleQuote{父啊！天地的主宰，我稱謝祢，因為祢將這些事瞞住了智慧及明達的人，而啟示給小孩子。是的，父啊！因為祢這樣願意，在祢面前就是如此。我父將一切都交給了我；除了父，沒有人認識子；除了子和子所願意啟示的人，沒有人認識父。你們所有勞苦和負重擔的，都到我這裡來，我要使你們安息。你們背起我的軛，跟我學吧，因為我是良善心謙的。}祂——父將一切都交給祂，除了父沒有人認識祂，唯獨祂（以及祂所願意啟示的人）認識父——祂不是說\BibleQuote{跟我學}去創造世界或復活死人，而是\BibleQuote{因為我是良善心謙的。}哦，具有救恩的教導！哦，人類的導師和主，人類因驕傲的杯爵而承擔並傳遞了死亡，祂不願教導祂自己所不是的，不願命令祂自己所不做的。我以祢為我所開啟的信德之眼，看見祢，良善的\Person{耶稣}，如同在人類的集會中呼喊說：\BibleQuote{你們到我這裡來，跟我學吧。}我懇求祢，萬物藉着祢而造的聖子，以及同一位在萬物中成了受造的聖子，我們到祢這裡來，向祢學什麼呢？祂說：\BibleQuote{因為我是良善心謙的。}難道藏在祢內的一切智慧和知識的寶藏，就是為了這個，叫我們向祢學一件大事，即祢是\BibleQuote{良善心謙的}？成為小的是如此偉大，以至於若非由如此偉大的祢來實現，就完全無法學到嗎？確實如此。因為除非那不安的膨脹——靈魂因之對自己顯得偉大，當它在祢面前不健康時——被消散，靈魂就找不到安息。\par

36. 让他们听到祢，让他们来到祢前，让他们向祢学习良善心谦，那些寻求祢的慈爱与真理的人，藉着为祢、而非为自己而生活。让那劳苦和负重担的人听到这话，他被自己的重担压得不敢举目望天，那个捶着胸膛、从远处走近的罪人。让他听到，那个百夫长，不配祢进入他的屋檐下。让他听到，匝凯，税吏长，将可诅咒之罪的所得四倍偿还。让她听到，那个城中的罪妇，在祢脚前泪流满溢，她离祢的脚步越远，就越多泪水。让他们听到，娼妓和税吏，他们在经师和法利塞人之前进入天国。让所有这类人听到，与他们的宴饮被当作把柄向祢提出指控，仿佛是由那些不需要医生的人提出的，而祢来不是召义人，而是召罪人悔改。所有这些人在归向祢时，都容易变得温良，在祢面前谦卑，想起自己最不义的生活和祢最宽厚的仁慈，因为\BibleQuote{罪恶在哪里增多，恩宠在那里也格外丰富}。\par

37. 但请看贞女的队伍，圣洁的男孩和女孩：这一群在祢的教会中接受培育；那里，他们从母亲的乳房就开始为祢发芽；为了祢的名，他们开口说话，祢的名，如同通过他们婴孩的乳汁，被注入并吮吸。他们中无人能说：\BibleQuote{我从前是亵渎者、迫害者和施暴者，但我蒙受了怜悯，因为我是在无知和不信中做的。}不仅如此，那祢没有命令、只是为那些愿意的人提出、使人可以抓住的，祢说：\BibleQuote{谁能领受，就让他领受吧}；他们抓住了，他们誓愿了，并且为了天国，不是因祢威胁，而是因祢劝勉，他们使自己成为阉者。向这些人呼喊，让他们听到祢，因为祢是\BibleQuote{良善心谦的}。不管他们多么伟大，让他们在一切事上如此谦卑自己，好能在祢面前获得恩宠。他们是义人，但难道他们能和祢一样，称不虔敬的人为义吗？他们是贞洁的，但他们的母亲在罪中孕育了他们。他们是圣的，但祢也是圣中之圣。他们是童贞者，但他们也不是由童贞女所生。他们在心神和肉体上完全贞洁，但他们不是成了血肉的圣言。然而，让他们学习，不是从那些祢赦免其罪的人，而是从祢自己，除免世罪的天主羔羊，因为祢是\BibleQuote{良善心谦的}。\par

38. 我并非差遣你，虔诚贞洁的灵魂，你没有给肉欲松缰直至允许婚姻，你没有放纵你即将离去的身体以至于生养继承者，你高举你尘世的肢体，漂浮着使它们习惯天上；我并非差遣你，为学习谦卑，到税吏和罪人那里去，他们尚且比骄傲的人先进天国；我不差遣你到他们那里：因为那些从不洁的深渊被释放出来的人，不配让无玷的童贞被差遣到他们那里去效法。我差遣你到天上的君王那里，到那创造人的那一位那里，也为人的缘故而在人中间受造；到那在世人中最美丽、却为了世人而被世人所轻蔑的那一位那里；到那统治永恒天使、却不嫌弃服事必死之人的那一位那里。总之，使他谦卑的不是不义，而是爱德；\BibleQuote{爱德不嫉妒，不自大，不求自己的益处}；因为\BibleQuote{基督也没有寻求自己的喜悦，而是正如经上论及祂所说的：\InnerQuote{辱骂你的人的辱骂都落在我身上。}}所以，去吧，到祂那里去，学习祂，因为祂是\BibleQuote{良善心谦的}。你不会到那个因不义的重压而不敢举目望天的人那里，而是到那位因爱德的分量从天降下的人那里。你不会到那位用眼泪洗主脚、寻求赦免重罪的人那里，而是到那位赦免一切罪、并洗自己门徒脚的那一位那里。我知道你童贞的尊贵；我没有向你提出谦卑地控告自己过错的税吏作为效法的榜样；但我惧怕那个骄傲地夸耀自己功绩的法利塞人。我不说：要像她，关于她曾说过：\BibleQuote{她的许多罪得了赦免，因为她爱得多}；但我恐怕你认为自己得赦免的少，因而爱得也少。\par

39. 我说，我大大为你们感到恐惧，惟恐当你们夸耀要跟随羔羊无论祂往何处去时，你们却因膨胀的骄傲而无法跟随祂走狭窄的道路。童贞的灵魂啊，这对你是好的：你既是童贞，因此在你的心中完全保持着你已重生，在你的肉身中保持着你已出生，你却仍怀有对主的敬畏，并产生救恩的精神。\Quote{恐惧}，诚然，\Quote{在爱德中没有恐惧，但完美的爱德}，正如经上记载，\Quote{除去恐惧}：但这是对人的恐惧，而非对天主的恐惧；对暂时祸患的恐惧，而非对最后审判的恐惧。\Quote{不要心高气傲，反要恐惧。}要爱天主的良善；要畏惧祂的严厉：两者都不允许你骄傲。因为借着爱，你恐惧，唯恐严重得罪那位被爱且爱你的主。因为有什么比因骄傲而得罪祂更严重的冒犯呢？祂为了你的缘故，曾使骄傲者不悦。在哪里更应有那\Quote{直到永远常存的纯洁敬畏}，岂不就是在你身上，你毫不思念世上的事，如何取悦婚姻伴侣，却思念主的事，如何取悦主？那种恐惧不在爱德中，但这种纯洁的敬畏不离弃爱德。若你不爱，你害怕灭亡；若你爱，你害怕不悦。那种恐惧，爱德将它驱逐；这种恐惧，爱德与它一同内住。宗徒保禄也说：\Quote{因为我们没有再次接受为恐惧作奴隶的精神，而是接受了义子的精神，我们借着它呼喊\InnerQuote{阿爸，父啊}}我相信他是指那种在旧约中赐下的恐惧，免得失去暂时的事物，就是天主曾应许给那些尚不在恩宠下作儿子、而是仍在法律下作奴隶的人的。也有对永恒之火的恐惧，为了避免它而侍奉天主，这肯定还不是完美的爱德。因为对奖赏的渴望是一回事，对惩罚的恐惧是另一回事。它们是不同的说法：\Quote{我往哪里去躲避祢的神？我往哪里逃，躲避祢的面？}和\Quote{我曾寻求主的一件事，我要寻求：就是一生一世住在主的殿中，瞻仰主的喜乐，并在祂的殿中得保护。}和\Quote{不要转脸不顾我。}和\Quote{我的灵魂渴慕且切望主的院宇。}那些话，让那不敢举目望天的人，以及那用眼泪洗祂的脚以求获得她重罪赦免的女人，去拥有吧；但你所拥有的，是那些挂念主的事，要在身体和灵魂上都成为圣洁的人。与那些话相伴的是有折磨的恐惧，完美的爱德将其逐出；但与这些话相伴的是对主的纯洁敬畏，它直到永远常存。对这两种人都必须说：\Quote{不要心高气傲，反要恐惧；}这样，人既不因自己的罪而自辩，也不因自以为义而自高。因为保禄自己，他说\Quote{你们没有再次接受为恐惧作奴隶的精神}，然而，作为爱德的同伴的恐惧，又说：\Quote{我在你们那里，又恐惧，又甚战兢。}而那句我曾提及的话，即接上的野橄榄枝不可向那折下的橄榄树枝骄傲，他自己也使用了，说：\Quote{不要心高气傲，反要恐惧；}他自己劝告基督的所有肢体，一般地说：\Quote{以恐惧战兢作成你们得救的工夫；因为是天主在你们心里运行，叫你们立志行事，为成就祂的美意。}免得那所记载的\Quote{当以敬畏侍奉主，以战兢喜乐地归向祂}似乎只属于旧约。\par

40. 圣身体，即教会，有哪些肢体比那些宣认童贞圣洁的人更应该谨慎，使圣神能安息在他们身上？但祂在哪里找不到祂自己的位置，又怎能安息呢？除了谦卑的心，祂能充满它，而非远离它；能提升它，而非压抑它？而经上说得最清楚：\Quote{我安息在谁身上？在那谦卑、安静、并对我的话战栗的人身上。}你已经正义地生活，你已经虔诚地生活，你贞洁、圣洁地生活，拥有童贞的纯洁；然而，你仍生活在此世，你听到\Quote{什么？人的生命在世上岂不是一个考验吗？}难道这不是使你从过分自信的傲慢中退却\Quote{祸哉，世界因了绊脚石！}难道你不战栗，惟恐你被列入那些\Quote{由于罪恶增多，许多人的爱情冷淡了}的人当中？当你听到\Quote{所以，那自以为站得稳的，要小心，免得跌倒。}你岂不捶胸？在这些神圣的警告和人类的危险之中，我们尚且很难说服圣洁的童贞女谦卑吗？\par

41. 或者，我们真的相信，天主容许许多男人和女人将要跌倒，与你们的修会人数混在一起，是为了别的什么原因，而不是为了借着这些人的跌倒，使你们的敬畏得以增加，以压制骄傲？而天主如此憎恨骄傲，以致于为了反抗这唯一的事，至高者竟自谦卑？除非，事实上，你们因此反而更少恐惧，更加自大，以至于少爱祂——祂曾如此爱你们，以至于为你们舍弃了自己——因为祂赦免了你们很少的罪，你们从小就度着虔诚、热心、贞洁、童贞未损的生活。仿佛你们真不应当以更炽热的爱去爱祂，祂没有让你们陷入那些罪恶，而这些罪恶在罪人归向祂时祂都赦免了。或者，那个法利塞人，认为自己被赦免的很少，因此爱得少，难道他不是因为这种错误而眼瞎，是因为他\Quote{不知道天主的正义，企图建立自己的正义，而没有服从天主的正义}吗？但你们，是特选的民族，在特选中更是特选的，童贞的合唱队跟随羔羊，连你们也\BibleQuote{藉着恩宠，藉着信德得救了，这不是出于你们自己，而是天主的恩赐：不是出于作为，免得有人自夸。因为我们是他的化工，在耶稣基督内受造，为行善事，就是天主所预备的，叫我们行在其中。}\BibleRef{弗 2:8-10}那么，你们越是因祂的恩赐而装饰，就越少爱祂吗？愿祂亲自阻止如此可怕的疯狂！因此，既然真理说了真理，那被赦免少的人，爱得少；那么，为了你们能以炽热的爱情去爱祂——你们从婚姻的束缚中被释放，可以自由地去爱——，你们要把一切因祂的引导而未曾犯的罪恶，都当作完全被赦免了。因为\BibleQuote{我的眼目常仰望上主，因为祂必将我的脚从网中拉出。}\BibleRef{咏 25:15}并且\BibleQuote{若不是上主看守城池，看守的人徒然警醒。}\BibleRef{咏 127:1}关于节德本身，宗徒说：\BibleQuote{我愿众人都像我一样；但每人都有天主所赐的恩宠，有人这样，有人那样。}\BibleRef{格前 7:7}那么，谁赐予这些恩赐呢？谁按祂的意愿把各自的恩赐分给每个人呢？当然是天主，祂没有不义，因此祂以何种公平使一些人这样，另一些人那样，人要知道要么不可能，要么十分困难；但祂是按公平而做的，这是不可怀疑的。\BibleQuote{你有什么不是领受的呢？}\BibleRef{格前 4:7}那么，你从祂那里领受得更多，却又因何等乖戾而少爱祂呢？\par

42. 所以，为穿上谦卑，首先应思想这一点：天主的童贞女不要以为自己之所以如此，是出于自己，而不当作是这最好的恩赐从上头、从光明之父降下来的，在他内没有变化或转动的阴影。因为这样她就以为赦免给她的很少，以致她爱得少，并且由于不知道天主的正义，而企图建立自己的正义，就不服从天主的正义。那个西满就犯了这错误，他被那女人超越了，那女人许多罪得了赦免，因为她爱得多。但她将有更谨慎、更正确的思想：我们应将所有罪恶看作就如已被赦免一样，因为天主保守我们不去犯这些罪。见证就是圣经中虔诚祈祷的词句，它们表明，那些受天主命令的事，若非借着命令者的恩赐和帮助，就不会成就。因为如果我们不靠祂恩宠的帮助就能做到，那么祈求这些事就是虚假的。有什么比遵守天主诫命的服从更普遍、更重要的命令呢？然而我们发现连这也被祈求。他说：\BibleQuote{你颁发了你的命令，叫他们严格地遵守。}随后说：\BibleQuote{但愿我的行径坚定，常遵守你的法令，我若注目于你的所有诫命，便不会蒙羞。}\BibleRef{咏 119:4-6}他前面记下天主所命令的，希望自己能完成。这确实是为了不犯罪；但如果犯了罪，命令就是要悔改；免得犯罪的人因辩护和推诿罪而死于骄傲，他因不愿自己所行的因悔改而消亡。这件事也向天主祈求，好使人明白，若不借着所祈求的那位的恩赐，就不能成就。\BibleQuote{上主，求你在我的口边派一守卫，防守我的嘴唇。莫使我的心倾向恶毒，与人同做奸犯科。}\BibleRef{咏 141:3-4}所以，如果遵守祂诫命的服从和使我们不为罪辩护的悔改，都是所愿所求的，那么显然，当这些成就时，是因祂的恩赐而拥有，因祂的帮助而实现。然而，关于服从，更明确地说：\BibleQuote{上主稳定人的脚步，喜悦人的道路。}\BibleRef{咏 37:23}关于悔改，宗徒说：\BibleQuote{或许天主要赐他们悔改。}\BibleRef{弟后 2:25}\par

43. 关于节德本身，岂非最清楚地说：\BibleQuote{我明知，除非天主赐予，没有人能节德；我知道这是谁的恩赐，这本身也是智慧的一部分？} 但也许节德是天主的恩赐，而智慧却是人自己赐予自己的，借此理解那恩赐不是自己的，而是来自天主的。是的，\BibleQuote{上主使盲人明智}；以及，\BibleQuote{上主的诫命是信实的，它赐智慧给幼稚的人}；以及，\BibleQuote{如果谁缺乏智慧，就该向天主祈求，祂厚赐众人而不责斥，必会赐给他}。但是贞女应当智慧，好使她们的灯不熄灭。如何\BibleQuote{智慧}？除非\BibleQuote{不心高气傲，却与卑微的人交往}。因为智慧本身曾对人说：\BibleQuote{看，虔诚就是智慧！} 因此，如果你没有什么是没有领受过的，\BibleQuote{不要心高气傲，倒要战战兢兢}。不要爱得少，仿佛那赦免你少的是祂；反而要多爱那赐予你多的主。因为如果那被赐予不必偿还的人爱，那么被赐予拥有的人更应当爱。因为凡是从起初就保持贞洁的，都由祂掌管；凡是从不洁变为贞洁的，都由祂改正；凡是一直不洁到底的，都被祂舍弃。但祂能通过隐秘的旨意这样做，却不能通过不义；也许为此原因它隐藏着，好使人更加敬畏，更少骄傲。\par

44. 再者，人既然知道靠天主的恩宠他才是他所是的，就不该落入另一个骄傲的陷阱，即为了天主的恩宠本身而高抬自己，从而鄙视其余的人。正是由于这个过失，那个法利塞人既为所拥有的美善感谢天主，却又在承认自己罪过的税吏之上自夸。因此，一个贞女应该做什么？她应该想什么，才不至于在那些没有这如此大恩赐的人（无论男女）之上自夸？因为她不应假装谦逊，而应彰显谦逊；因为假装的谦逊是更大的骄傲。因此，圣经愿意显示谦逊应当是真实的，在说了\BibleQuote{你越伟大，在一切事上越当谦下}之后，紧接着说\BibleQuote{这样，你必在天主面前获得恩宠}；确实，在那里人不能欺骗地谦下。\par

45. 那么我们该说什么？有没有一种思想，天主的贞女可以真实地拥有，因而她不敢将自己置于一位忠信的女子之前——不仅是寡妇，甚至也包括已婚者？我不是说一个被弃绝的贞女；因为谁不知道一个服从的女子应当置于一个不服从的贞女之前？但当两者都服从天主的命令时，她会如此战栗，不敢将圣洁的童贞置于贞洁的婚姻之上、将节德置于婚姻生活之上、将百倍的果实置于三十倍之前吗？不，让她不要怀疑将此物置于彼物之上；然而，这位或那位敬畏天主的贞女，不敢将自己置于那位或这位敬畏天主的女子之前；否则她就不会谦卑，而\BibleQuote{天主相反骄傲的人！} 那么，她心里该想什么？实际上是天主的隐藏恩赐，除了试探的询问，没有人能知道，即使在自己身上。因为，暂且不提其他，一个贞女虽然挂虑主的事，为悦乐主，但怎能知道，由于她自己都不知道的心灵弱点，她是否还尚未准备好殉道，而那个她乐于置身其前的女子，却可能已经能喝主羞辱之杯——就是主在门徒贪求高位时放在他们面前要他们先喝的那杯？我说，她怎能知道，她自己还不是\Person{德克拉}，而那个女子已经是\Person{克里斯比纳}？当然，除非有眼前的试炼，否则这恩赐就无法得到证明。\par

46. 但这恩赐如此伟大，以致有人认为它是百倍的果实。因为教会的权威给出了非常显著的见证：在祭台圣事中，信友知道在哪里纪念殉道者，在哪里纪念已故的圣洁修女。至于这果实差异的含义，让那些比我们更了解这些事的人去判断吧：是童贞的生活为百倍果实，守寡为六十倍，婚姻为三十倍；还是百倍果实归于殉道，六十倍归于节德，三十倍归于婚姻；或者，童贞加上殉道才达到百倍，而单独的童贞为六十倍，但婚姻中的人若成为殉道者，则从三十倍达到六十倍；或者，在我看来更可能的是，因为天主恩宠的恩赐众多，且一个比另一个更大更好，因此宗徒说：\BibleQuote{但你们应追求那更大的恩赐}（\BibleRef{格前 12:31}），我们应理解这些恩赐的数量之多，无法被归入那些不同的类别。首先，我们不将寡居的节德视为没有果实，也不将其与婚姻爱德的功绩相提并论，或与童贞的光荣相等；也不认为殉道者的冠冕，无论是已有的心神习惯（即使缺乏考验的证明），还是实际经受苦难的考验，加于任何一种贞洁之上而没有任何果实的增加。其次，当我们认定许多男女如此持守童贞，却仍不做主所说的那些事：\BibleQuote{你若愿意是成全的，去！变卖你所有的，施舍给穷人，你必有宝藏在天上，然后来跟随我}（\BibleRef{玛 19:21}），并且不敢与那些共同生活的人联合，在他们中间\BibleQuote{没有人说自己的财物归自己，一切都归公用}（\BibleRef{宗 4:32}）；难道我们认为当天主的贞女这样做时，她们的果实没有任何增加吗？或者虽然她们不这样做，她们就没有任何果实吗？因此，恩赐很多，有些比其他的更光明、更崇高。有时，一个人在较少的恩赐上结果实，但更好；另一个人在较低的恩赐上结果实，但更多。而在领受永恒荣耀时，它们彼此如何相等或区分，有谁敢断言呢？然而，显然这些差异很多，且更好的恩赐并非为现世，而是为永恒有益。但我认为，主有意提到三种果实差异，其余的留给了那些能理解的人。因为另一位福音作者只提到了百倍；难道我们因此认为他拒绝了或不知道其他两种吗？相反，他留下它们让人理解。\par

47. 但是，正如我开始所说的，无论百倍的果实是献于天主的童贞，还是我们以其他方式理解这果实的间隔（或是我们已提及的，或是我们未提及的），然而，我相信没有人会敢将童贞置于殉道之上，也没有人会怀疑，如果缺乏考验的试验，这后一种恩赐是隐藏的。因此，贞女有一个值得思考的题材，可能有利于她保持谦卑，以免她违背那超越一切恩赐的爱德，没有这爱德，她所拥有的其他任何恩赐，无论多少，无论大小，都算不得什么。我说，她有一个值得思考的题材：以免她自高自大，嫉妒争竞；确实，她如此宣称童贞之善远大于婚姻之善，却不知道这个或那个已婚妇女是否已经能够为基督受苦，而她自己却还不能，她在此受到宽免，她的软弱没有受到考验的质问。因为宗徒说：\BibleQuote{天主是忠信的，他决不许你们受那超过你们能力的试探；他如加给人试探，也必开一条出路，叫你们能够承担}（\BibleRef{格前 10:13}）。因此，或许那些持守婚姻生活（在其类中值得赞美）的男女，已经能够抵挡不义的敌人，即使肠破血流也要争战；但这些自幼节德、为天国自阉的男女，仍然未能如此忍受，无论是为正义还是为贞洁本身。因为，为了真理和神圣目的，不服从一个劝说和谄媚的人，是一回事；而不屈服于一个拷打和击打的人，是另一回事。这些隐藏在灵魂的能力和力量中，通过考验才显明，通过实际经历才显现。因此，为了每个人不因自己能清楚做到的事而自高自大，让他谦卑地考虑：他不知道或许存在一些更卓越的事，他做不到；但有些人，既不拥有也不宣称他合法自觉的事，却能做他自己不能做的事。这样，靠着真实的谦卑而非虚假的，就能遵守：\BibleQuote{彼此以恭敬相待}（\BibleRef{罗 12:10}）和\BibleQuote{存心谦卑，各人看别人比自己强}（\BibleRef{斐 2:3}）。\par

48. 关于对罪的极其谨慎和警醒，我现在该说什么呢？\BibleQuote{谁能夸口说自己有贞洁的心？谁能夸口说自己无罪可染？}圣洁的童贞确实在母胎中便已无玷；但他说：\BibleQuote{在祢眼前，没有一人是洁净的，连那在世上一日的婴孩也不例外。}在信德中，也持守一种无玷的童贞，借此教会如同一贞洁的童女与一位新郎结合。但那位独一新郎教导的不仅是那些身心贞洁的信徒，而是所有基督徒，从属灵到属肉体的，从宗徒到最后忏悔的罪人，如同从天界的高处到它的边界，都要祈祷；并且在祈祷中，祂告诫他们要说：\BibleQuote{请宽恕我们的罪债，如同我们也宽恕我们的债务人一样。}在此，我们藉由所寻求的，祂也表明我们必须记得我们是什么人。因为祂并非要我们为那些我们相信在圣洗中通过祂的和好已得赦免的过去一生的罪债祈祷，说：\BibleQuote{请宽恕我们的罪债，如同我们也宽恕我们的债务人一样。}否则这祈祷更适合慕道者在接受圣洗之前祈祷；然而，这是受过洗的人祈祷的内容，包括统治者与人民、牧者与羊群；这便充分表明，在整个人生（即是一场考验）中，没有人能夸耀自己全然无罪。\par

49. 为此，天主的贞女诚然无瑕，\BibleQuote{随从羔羊往哪里去，}既成全了罪的洁净，又保持了童贞——一旦失去便不能复得。但因为同一部\BibleBook{}{默示录}（其中如此事向某人启示）也称赞她们，说\BibleQuote{在他们口中找不出谎言，}让她们也记得在此事上要真实，不敢说自己无罪。因为那位看见此事的同一\Person{若望}说了这话：\BibleQuote{若我们说我们没有罪，便是自欺，真理不在我们内；但若我们告明我们的过犯，祂是忠信正义的，必赦免我们的罪，并洗净我们的一切不义。若我们说我们没有犯过罪，我们便是以祂为说谎者，祂的话就不在我们内。}这话当然不是针对某些人，而是针对所有基督徒，贞女也该在其中认出自己。因为如此，她们将如\BibleBook{}{默示录}中所显示的那样，成为无谎言的人。以此方式，只要在天界的高处尚未达到完美，谦卑中的告明便使她们无过。\par

50. 然而，恐怕有人借这段话之故，怀着致命的安逸去犯罪，让自己被冲昏头脑，仿佛其罪过很快就能通过轻易的告明而被涂抹，他（若望）立刻补充说：\BibleQuote{我的孩子们，我给你们写这些，是要你们不犯罪；若有人犯了罪，我们在父那里有辩护者，就是正义的耶稣基督，祂自己就是我们的赎罪祭。}因此，谁也不要离开罪又准备返回，也不要与不义缔结这样的联盟契约，以至于更喜欢告明罪而不是躲避罪。然而，即使对那些忙碌警醒不犯罪的人，由于人性的软弱，仍会有一些罪悄悄潜入，无论多么微小、稀少，也并非全无；这些罪若被骄傲加重加重，就变得重大而严重；但如果藉着虔诚的谦卑在我们在天上的大司祭面前被摧毁，便极容易得到洁净。\par

51. 但我并不与那些宣称人能在此生无罪生活的人争辩；我不争辩，也不反对。因为我们或许是从自己的可怜境遇去度量伟人，与自己比较，不明所以。我知道一事：那些伟人，如同我们所不是的、我们尚未证明的，他们越是伟大，就越在一切事上谦卑自己，以便在天主面前获得恩宠。因为无论他们多么伟大，\BibleQuote{没有仆人大于主人，也没有门徒大于师傅。}而祂确实是主，祂说：\BibleQuote{一切都由我父交付给我了。}祂是师傅，祂说：\BibleQuote{你们所有劳苦的人，到我这里来，跟我学习。}而我们学习什么呢？祂说：\BibleQuote{因为我良善心谦。}\par

52. 有人会说：现在写的不是童贞，而是谦卑。仿佛我们所着意陈述的，不是任何童贞，而是那属于天主的童贞似的。这善愈是伟大，我就愈是为它担忧，怕它因骄傲这个盗贼而丢失。因此，除了赐予童贞的天主本身，没有人能守护童贞之善；而天主是爱德。所以童贞的守护者是爱德；而这守护者的处所是谦卑。祂确实住在那里，祂曾说：\BibleQuote{我的精神安息在谦卑、安静、且畏惧我言语的人身上。}因此，我若希望我所赞扬的善得到更稳妥的守护，而留意为守护者预备一个处所，这哪里背离了我的目的呢？因为我自信地说，也不怕那些我以关爱劝诫要与我一同学会惧怕的人向我发怒。谦卑的已婚者，虽然不能像羔羊那样无论祂往哪里去都跟随，但尽可能跟随，比起骄傲的童贞女更容易跟随羔羊。因为一个人若不愿接近祂，怎能跟随祂？若不来学祂的\BibleQuote{良善心谦}，又怎能接近祂？因此，羔羊带领那些跟随祂的人，无论祂往哪里去，首先必须在他们中间找到一个可以安放祂头颅的地方。因为有一个骄傲而狡猾的人曾对祂说：\BibleQuote{主，祢无论往哪里去，我都要跟随祢。}祂回答说：\BibleQuote{狐狸有洞，天空的飞鸟有巢，但人子却没有枕头的地方。}祂以狐狸之名斥责诡诈的狡猾，以飞鸟之名斥责膨胀的骄傲，在这些人中祂找不到虔诚的谦卑来安息。因此，那个答应要跟随祂的人，没有在任何地方跟随主，因为他不是只跟随一段路程，而是要随祂往任何地方去。\par

53. 因此，天主的贞女们，你们要这样做，要这样做：跟随羔羊，无论祂往哪里去。但你们先要来到祂面前，你们将要跟随的那一位，并学习祂的良善心谦。你们要谦卑地来到谦卑者面前，如果你们爱祂；不要离开祂，免得跌倒。因为害怕离开祂的人祈求说：\BibleQuote{不要让骄傲的脚临到我。}你们要以谦卑的脚行走在高处的路上；祂亲自在谦卑中高举那些跟随的人，祂不觉得降卑到卑微者中间是麻烦。你们要把祂的恩赐交托给祂保管，\BibleQuote{向祂守护你们的力量。}凡是因祂的守护而你们没有犯的恶，都算作祂赦免了你们；免得你们以为欠祂的少，就爱得少，并以毁灭性的夸耀轻视那些捶胸的税吏。你们要谨慎保守那经受过考验的力量，免得因你们能够承受某些事而自高自大；至于那未经考验的，你们要祈祷，免得你们受试探超过你们所能承受的。要想到有些人在暗处比你们更优秀，而你们在明处比他们更好。当你们善意地相信他人那些你们也许不知道的善事时，你们自己所知道的善事不会因比较而减少，反而因爱德而增强；而那些也许尚欠缺的，也会因更谦卑地渴望而更容易获得。愿你们中坚持的人成为你们的榜样；愿那些跌倒的人增加你们的恐惧。你们要爱那榜样，好能效法它；为那跌倒者哀悼，免得你们自高自大。你们不要建立自己的义，要顺服那使你们成义的天主。宽恕他人的罪过，为自己的罪过祈祷；借着警醒躲避未来的罪过，借着告解涂抹过去的罪过。\par

54. 看，你们已经如此，以致在其余的行为中也与你们所宣发并持守的童贞相符。看，你们不仅不犯杀人、魔鬼的祭献和憎恶、偷窃、抢劫、欺诈、伪誓、酗酒，以及一切纵欲和贪婪、仇恨、嫉妒、不敬、残忍；而且甚至连那些被视为或被认为是较轻的事，在你们中间也找不到，也不出现：不厚颜，不眼目游荡，不口无遮拦，不嬉笑轻浮，不下流玩笑，不举止失当，不昂首阔步或松散行走；你们已经不以恶报恶，不以咒骂还咒骂；最后，你们已经达到爱德的尺度，甚至为弟兄舍命。看，你们已经如此，因为你们也应当如此。这些与童贞结合，在人间展现了天使般的生活，在天上展现了通往天堂的道路。但是，不论你们多么伟大，你们中凡如此伟大的，\BibleQuote{在一切事上越当谦卑自己，好能在天主面前获得恩宠，}免得祂作为骄傲者抵挡你们，免得祂因你们自高而降卑你们，免得祂因你们膨胀而引导你们经过窄路：尽管不必焦虑，哪里有爱德燃烧，哪里就不会缺少谦卑。\par

55. 因此，如果你们藐视人子的婚姻（从婚姻而生的人子），你们要以全心爱那在世人中俊美的那一位；你们有空闲，你们的心不受婚姻的束缚。注视你们爱人的美丽：思想祂与圣父同等，却也服从祂的母亲；在天上治理，在地上服事；创造万有，在万有中受造。骄傲者所嘲笑的那一点，你们要注视它多么美丽：以内在的眼目注视被钉者的伤口，复活着伤痕，死者的血，相信者的代价，救赎者的收益。思量这些有何等大的价值，用爱德的天平称量；然后把你们原本要花费在婚姻上的爱情，全部偿还给祂。\par

56. 祂寻找你们内在的美丽，这很好，因为祂赐给了你们成为天主女儿的能力；祂向你们要求的不是美丽的肉身，而是美好的品行，由此也能驾驭肉身。祂不是那种可以被任何人用谎言诽谤你们、并使祂因妒忌而发怒的。你们看，你们爱祂是多么安全，你们不怕因虚假的猜疑而得罪祂。丈夫和妻子彼此相爱，因为他们彼此看见；而对于看不见的事，他们彼此担忧；他们在可见之事上没有确定的喜乐，而在隐藏之事上通常猜疑不存在的事。你们在祂内，你们用肉眼看不见祂，却凭信德注视祂，既没有真实的过错可以责备，也不怕或许因虚假而得罪祂。所以，如果你们应当对丈夫付出大爱，那么对为了祂的缘故你们不愿有丈夫的那一位，你们应当爱得多么深呢？让祂安放在你们的整个心中，祂为了你们曾被钉在十字架上；让祂占有你们灵魂中一切你们不愿被婚姻占据的部分。你们不可少爱祂，为了祂的缘故你们甚至没有爱那本合法之事。既然爱那良善心谦的主，我就不为你们担忧骄傲。\par

57. 因此，按我们微薄的能力，我们已充分谈论了圣洁（藉此你们被恰当地称为\Quote{\OriginalTerm{献身者}{sanctimoniales}}）和谦卑（藉此你们所负的任何大名得以保持）。但让那三位青年更有价值地劝诫你们关于我们这篇小作品，他们以心中炽热的爱所爱的那一位在火中赐予了他们凉风；他们劝诫你们，在言辞上确实简短得多，但在权威的分量上却大得多，体现在那首赞歌中，天主藉他们受尊崇。因为他们将谦卑与圣洁结合于赞美天主的人身上，极其清楚地教导：每个人，无论他做出何等神圣的宣认，就越要谨慎，以免被骄傲欺骗。因此，你们也要赞美祂，祂赐予你们，在这世界的火焰中，虽然你们未缔结婚姻，却未被焚烧；并为我们祈祷：\Quote{请你们祝福主，你们这些心地圣洁谦卑的人；请向主咏唱赞歌，永远颂扬祂超乎万有。}\par

\SourceRecord{Translated by C.L. Cornish.；Nicene and Post-Nicene Fathers, First Series；Vol. 3.；Edited by Philip Schaff.；Buffalo, NY: Christian Literature Publishing Co.,；1887.；<http://www.newadvent.org/fathers/1310.htm>.}
